\documentclass[a4paper,10pt]{article}
\usepackage[margin=0.5in,nofoot]{geometry}
\usepackage{fontawesome5}
\usepackage{hyperref}
\usepackage{titlesec}
\usepackage{xcolor}
\usepackage[utf8]{inputenc}

\hypersetup{
    colorlinks=true,
    linkcolor=black,
    filecolor=black,
    urlcolor=black,
    citecolor=black
}

\titleformat{\section}{\large\bfseries}{\thesection}{1em}{}[\titlerule]
\titlespacing*{\section}{0pt}{*1}{*1}

\newcommand{\entry}[4]{
  \noindent\textbf{#1} \hfill #2 \\
  \noindent\textit{#3} \hfill \textit{#4} \\
  \vspace{2pt}
}

\begin{document}

\pagenumbering{gobble}

\noindent
\begin{minipage}[t]{0.5\textwidth}
\textbf{\Huge Gabriel Canela}

\vspace{0.4em}

\end{minipage}%
\begin{minipage}[t]{0.5\textwidth}
\raggedleft

Campinas - SP

\href{tel:+5519996831296}{\faPhone \space +55 (19) 99683-1296}

\href{mailto:gabriel.canela.ti@gmail.com}{\faEnvelope \space gabriel.canela.ti@gmail.com}

\vspace{0.2em}

\href{https://github.com/Canela-san}{\faGithub \space GitHub - Canela-san} \quad
\href{https://www.linkedin.com/in/gabriel-canela-ti/}{\faLinkedin \space LinkedIn}
\end{minipage}

\vspace{1em}

\begin{center}
    \textbf{\Large Eng. de Controle e Automação | Mecatrônica e Modelagem}
\end{center}

\vspace{0.5em}

\section*{Resumo Profissional}

Estudante de Engenharia de Controle e Automação pela Unicamp, com formação técnica dupla em Informática e Informática para Internet (ETEC) e perfil focado no desenvolvimento de sistemas de precisão. Atualmente atuo como Pesquisador (Bolsista FAPESP) no desenvolvimento de instrumentação embarcada para redes elétricas, combinando modelagem de sistemas de controle, integração mecatrônica e processamento digital de sinais.

Possuo experiência prática em laboratório de pesquisa (LabREI/Unicamp), atuando na montagem e validação de conversores de potência. Proficiente em programação de baixo nível (C, Assembly RISC-V/PRU) e alto nível (Python, scripts analíticos), além de domínio em ferramentas de simulação e projeto de hardware (MATLAB/Simulink, Altium Designer, LTspice).


\section*{Experiência}

\entry{FEEC Unicamp (Projeto FAPESP)}{\faCalendar \space 2025 -- Atual}{Pesquisador de Iniciação Científica}{}
\vspace{-1.6em}
\begin{itemize}
\setlength\itemsep{0em}
\item Desenvolvimento de um sistema embarcado de medição de alta performance para identificação de supraharmônicos em redes elétricas.
\item Programação de firmware em C e Assembly no subsistema de tempo real (PRU-ICSS) da BeagleBone Black, controlando o barramento SPI para aquisição de dados determinística.
\item Execução de engenharia reversa e depuração física em placas de circuito impresso (PCB) customizadas de aquisição de sinal.
\item Reestruturação de software embarcado em ambiente Linux, com reconfiguração de Device Tree e integração de processos de persistência de dados.
\item Elaboração de setups de testes e condução de validações experimentais baseadas em processamento digital de sinais (FFT).
\end{itemize}

\entry{LabREI (Lab. de Redes Elétricas Inteligentes) - FEEC Unicamp}{\faCalendar \space 2022 -- 2024}{Técnico de Laboratório}{}
\vspace{-1.6em}
\begin{itemize}
\setlength\itemsep{0em}
\item Apoio técnico a pesquisadores de pós-graduação na elaboração de setups de testes e validação experimental de projetos de eletrônica de potência.
\item Integração de sistemas mecânicos e eletrônicos: cotação e escolha de componentes, usinagem/furação de painéis metálicos e execução das conexões elétricas de um conversor CA/CC bidirecional de 4 pernas de alta potência.
\item Administração da infraestrutura de TI do laboratório, gerenciando servidor local baseado em Linux Ubuntu, site institucional e atualização de certificados.
\item Treinamento de usuários na correta operação de equipamentos de medição e manutenção de hardware/software dos computadores de pesquisa.
\end{itemize}

\section*{Projetos e Destaques}

\entry{Controle de Velocidade de Motor DC}{[S1/2026]}{Projeto Acadêmico -- Laboratório de Controle (Unicamp)}{}
\vspace{-1.6em}
\begin{itemize}
\setlength\itemsep{0em}
\item Projeto e sintonia de um controlador PID para controle de velocidade de um motor DC, otimizado para atingir e frear na velocidade de referência no menor tempo possível.
\item Modelagem da planta e ajuste dos ganhos do controlador pelo método do lugar das raízes (\textit{root locus}), com simulação da resposta em malha fechada em MATLAB/Simulink.
\item Análise de tempo de acomodação, sobressinal (\textit{overshoot}) e estabilidade do sistema controlado.
\end{itemize}

\entry{SH-Analyzer: Analisador de Supraharmônicos}{2025 -- Atual}{Iniciação Científica (Bolsista FAPESP)}{}
\vspace{-1.6em}
\begin{itemize}
\setlength\itemsep{0em}
\item Desenvolvimento de um sistema híbrido de instrumentação de alto desempenho (BeagleBone Black + PCB customizada) para aquisição determinística de sinais em redes elétricas.
\item Programação de firmware em Assembly nativo para o coprocessador de tempo real PRU-ICSS (PRU 0), otimizando a comunicação via barramento SPI com o ADC ADS8688 de 16 bits.
\item Implementação de sincronização ARM-PRU via memória compartilhada interna e rotinas em Python (Numpy/Scipy) para processamento espectral via FFT e calibração de dados brutos.
\end{itemize}

\entry{Custom Processor from Scratch with MOSFETs}{2025 -- Atual}{Projeto de Microarquitetura e Hardware}{}
\vspace{-1.6em}
\begin{itemize}
\setlength\itemsep{0em}
\item Concepção de arquitetura computacional completa a partir do nível de componentes físicos, mapeando diagramas de blocos de ULA, registradores e unidade de controle.
\item Modelagem e simulação completa de circuitos lógicos digitais no software \textit{Digital}, realizando a transição física e validação de portas lógicas puramente baseadas em transistores MOSFET.
\item Desenvolvimento de um conjunto de instruções (ISA) e montador Assembly customizado para validação física da microarquitetura através da emulação de periféricos de display.
\end{itemize}

\entry{Nuvem-Privada: Infraestrutura NAS Self-Hosted}{2026}{Projeto de Automação de Infraestrutura de TI}{}
\vspace{-1.6em}
\begin{itemize}
\setlength\itemsep{0em}
\item Orquestração e containerização de um ecossistema de nuvem privada utilizando Docker e Docker Compose, integrando Nextcloud, bancos de dados MariaDB e cache Redis com segurança baseada em health checks.
\item Automação de rotinas de manutenção em background via Shell Script, implementando pipelines de sincronização de arquivos físicos, limpeza de banco de dados e backups consistentes a frio.
\item Provisionamento seguro de rede através de VPN mesh (Tailscale) para acesso remoto criptografado sem exposição direta de portas lógicas do servidor à internet.
\end{itemize}

\section*{Educação}

\entry{Unicamp Campinas}{\faCalendar \space 2022 -- 2028 (previsto)}{Engenharia de Controle e Automação}{}

\entry{ETEC}{\faCalendar \space 2018 -- 2019}{Técnico em Informática}{}

\entry{ETEC}{\faCalendar \space 2017 -- 2019}{Técnico em Informática para Internet}{}

\section*{Habilidades e Competências}

\begin{itemize}
\setlength\itemsep{0em}
\item \textbf{Controle, Automação e Modelagem:} Modelagem matemática e simulação de sistemas dinâmicos; projeto, sintonia e validação de controladores em malha fechada (PID, método do lugar das raízes) em MATLAB e Simulink; simulação de circuitos digitais e transitórios em Digital e LTspice.
\item \textbf{Sistemas Embarcados e Tempo Real:} Arquiteturas de microprocessadores e microcontroladores (ARM Cortex-A8, PRU-ICSS, RISC-V); programação de firmware em C e Assembly nativo; mapeamento de registradores, controle determinístico de jitter e comunicação via barramento SPI.
\item \textbf{Sistemas Operacionais Aplicados:} Desenvolvimento embarcado em ambiente Linux (Baseados em Debian); configuração avançada de multiplexação de pinos (Pinmux) e Árvore de Dispositivos (Device Tree Overlay); administração de servidores locais Linux Ubuntu e SSH/SFTP.
\item \textbf{Eletrônica e Instrumentação:} Projetos, modelagem e simulação de esquemáticos e placas de circuito impresso (PCB) multicamadas no Altium Designer; técnicas de isolamento galvânico (dados e potência), level shifting e condicionamento analógico de sinais; diagnóstico físico de circuitos com osciloscópio e multímetro.
\item \textbf{Programação e Ciência de Dados:} Python avançado para automação de scripts, engenharia de dados de séries temporais e pós-processamento analítico (Pandas, NumPy, SciPy); Processamento Digital de Sinais (DSP) focado em análise espectral via Transformada Rápida de Fourier (FFT) e filtragem digital; consultas estruturadas em SQL.
\item \textbf{DevOps e Ferramentas de Engenharia:} Versionamento e integração contínua (Git/GitHub); containerização e orquestração de microsserviços (Docker e Docker Compose) com automação em Shell Script (Bash); ambientes de desenvolvimento VS Code.
\end{itemize}

\section*{Cursos, Certificações e Idiomas}

\begin{itemize}
\setlength\itemsep{0em}
\item \textbf{Inglês -- Avançado (Nível Master Fluency):} Certificação de conclusão de ciclo avançado nível 11 — CCAA (2015 -- 2020).
\item \textbf{Japonês -- Básico (JLPT N5):} \textit{Japanese-Language Proficiency Test} — Certificado oficial de proficiência emitido pela Japan Foundation.
\item \textbf{Formação Avançada em TI:} Certificados de Qualificação Profissional em Desenvolvimento Web, Arquitetura de Redes e Banco de Dados integrados ao Centro Paula Souza (ETEC).
\end{itemize}

\end{document}